\documentclass[UTF8,10pt,letterpaper,fontset=fandol]{ctexart}
\usepackage[margin=0.9in]{geometry}
\usepackage{amsmath,amssymb,enumitem,needspace}
\usepackage[unicode,hidelinks]{hyperref}
\pagestyle{plain}
\raggedbottom
\setlength{\parindent}{0pt}
\setlength{\parskip}{1pt}
\setlength{\emergencystretch}{2em}
\linespread{1.06}
\AtBeginDocument{\setlength{\abovedisplayskip}{5pt}\setlength{\belowdisplayskip}{5pt}\setlength{\abovedisplayshortskip}{3pt}\setlength{\belowdisplayshortskip}{3pt}}
\setcounter{secnumdepth}{0}
\setlist[enumerate]{label=(\arabic*),leftmargin=2em,itemsep=0.2em,topsep=0.2em}
\ctexset{section={format=\Large\bfseries,beforeskip=1.1ex,afterskip=0.7ex}}
\newcommand{\C}{\mathbb C}
\newcommand{\R}{\mathbb R}
\newcommand{\F}{\mathbb F}
\newcommand{\ip}[2]{\langle #1,#2\rangle}
\DeclareMathOperator{\rank}{rank}
\DeclareMathOperator{\tr}{tr}
\DeclareMathOperator{\im}{im}
\begin{document}
\begin{center}
{\LARGE 作业 12}\\[0.5em]
2026 年 11 月 13 日\\[0.25em]
姓名：YuanHao Tang\quad 学号：3120265621
\end{center}
\section{第 1 题}
证明定理 5.10。设 $A\in\C^{m\times n}$ 的秩为 $r$。证明存在列满秩的 $B\in\C^{m\times r}$ 与行满秩的 $C\in\C^{r\times n}$ 满足 $A=BC$。给出构造下列可逆矩阵 $P,Q$ 的实际步骤，并验证所给两因子的秩：\par
\[
 PAQ=\begin{pmatrix}I_r&D\\0&0\end{pmatrix},\qquad
 B=P^{-1}\binom{I_r}{0},\qquad C=[\,I_r\ D\,]Q^{-1}.
\]
\par\Needspace{4\baselineskip}\textbf{解答}\par
行消元找主元；用列置换 $Q$ 将原矩阵的主元列移至前方：$AQ=[B_0\ N]$。因 $B_0$ 为列空间的基，$N=B_0D$。\par
行化简 $[B_0\mid I_m]\longrightarrow[\binom{I_r}{0}\mid P]$。于是 $PN=\binom D0$，从而\par
\[
 A=P^{-1}\binom{I_r}{0}[\,I_r\ D\,]Q^{-1}=BC.
\]
两内因子均含 $I_r$，$P,Q$ 可逆，故 $\rank B=\rank C=r$（$r=0$ 时取空因子）。\par
\section{第 2 题}
对下列矩阵作满秩分解。写出消去第一列的 $T_1$、消去第三列其余元素的 $T_2$、$P=T_2T_1$、$P^{-1}$ 和两个因子。\par
\[
 A=\begin{pmatrix}1&2&1&1&4\\2&4&3&0&8\\3&6&4&1&12\end{pmatrix}.
\]
\par\Needspace{4\baselineskip}\textbf{解答}\par
先作 $R_2\gets R_2-2R_1,\ R_3\gets R_3-3R_1$，再作 $R_1\gets R_1-R_2,\ R_3\gets R_3-R_2$：\par
\[
T_1=\begin{pmatrix}1&0&0\\-2&1&0\\-3&0&1\end{pmatrix},\quad
T_1A=\begin{pmatrix}1&2&1&1&4\\0&0&1&-2&0\\0&0&1&-2&0\end{pmatrix},\quad
T_2=\begin{pmatrix}1&-1&0\\0&1&0\\0&-1&1\end{pmatrix}.
\]
\[
 P=\begin{pmatrix}3&-1&0\\-2&1&0\\-1&-1&1\end{pmatrix},\quad
 P^{-1}=\begin{pmatrix}1&1&0\\2&3&0\\3&4&1\end{pmatrix},\quad
 PA=\begin{pmatrix}1&2&0&3&4\\0&0&1&-2&0\\0&0&0&0&0\end{pmatrix}.
\]
$B$ 取 $P^{-1}$ 的前两列，$C$ 取 $PA$ 的非零行：\par
\[
 \boxed{A=BC=\begin{pmatrix}1&1\\2&3\\3&4\end{pmatrix}
 \begin{pmatrix}1&2&0&3&4\\0&0&1&-2&0\end{pmatrix}} .
\]
\newpage
\section{第 3 题}
证明引理 5.1：对任意复矩阵 $A$，$\rank(AA^*)=\rank(A^*A)=\rank A$。\par
\par\Needspace{4\baselineskip}\textbf{解答}\par
对任意 $x\in\C^n$，\par
\[
 A^*Ax=0\ \Longrightarrow\ 0=x^*A^*Ax=\|Ax\|^2
 \ \Longrightarrow\ Ax=0,\qquad \ker(A^*A)=\ker A.
\]
秩–零度定理给出 $\rank(A^*A)=\rank A$。同理 $\rank(AA^*)=\rank A^*$；由 $\ker A^*=(\im A)^\perp$，$\rank A^*=\rank A$。\par
\section{第 4 题}
分别求原题 (a)、(d) 中矩阵的满秩分解。\par
\[
 A_1=\begin{pmatrix}2&5&8&11&14\\2&4&6&8&10\\4&9&14&19&24\end{pmatrix},\qquad
 A_4=\begin{pmatrix}-1&-1&-1&-1&-3&-1\\1&1&1&1&-1&-3\\3&3&3&3&1&-1\\3&3&3&3&1&-1\\3&3&3&3&1&-1\\3&3&3&3&1&-1\end{pmatrix}.
\]
\par\Needspace{4\baselineskip}\textbf{解答}\par
\textbf{(a)} 首两列的前两行子式为 $-2$。由 $a_k=a_1+(k-1)(a_2-a_1)$ 得\par
\[
 \boxed{A_1=\begin{pmatrix}2&5\\2&4\\4&9\end{pmatrix}
 \begin{pmatrix}1&0&-1&-2&-3\\0&1&2&3&4\end{pmatrix}} .
\]
右因子的前两列组成 $I_2$，两因子秩均为 2。\par
\textbf{(d)} 取原矩阵第 $1,5,6$ 列，其前三行子式为 $16$。前四列相同，故\par
\[
 \boxed{A_4=
 \begin{pmatrix}-1&-3&-1\\1&-1&-3\\3&1&-1\\3&1&-1\\3&1&-1\\3&1&-1\end{pmatrix}
 \begin{pmatrix}1&1&1&1&0&0\\0&0&0&0&1&0\\0&0&0&0&0&1\end{pmatrix}} .
\]
右因子第 $1,5,6$ 列组成 $I_3$，两因子秩均为 3。\par
\end{document}
